% TOP_PROBLEM: {"problemId": "problem.piltz-divisor-problem-for-k-fold-divisor-error-terms", "canonicalTitle": "Piltz Divisor Problem for k-Fold Divisor Error Terms", "releaseRank": 438, "primaryDomain": "number_theory_arithmetic_geometry", "primaryDomainLabel": "Number theory & arithmetic geometry", "displayStatus": "open", "releaseStatus": "open"}
% !TeX program = lualatex
% ATTEMPT_STATUS: unresolved
\documentclass[11pt]{article}
\usepackage[margin=25mm]{geometry}
\usepackage{fontspec}
\setmainfont{Latin Modern Roman}
\usepackage{amsmath,amssymb,mathtools}
\usepackage{unicode-math}
\setmathfont{Latin Modern Math}
\usepackage{xurl,hyperref}
\hypersetup{colorlinks=true,urlcolor=blue}
\setcounter{secnumdepth}{0}
\setlength{\parindent}{0pt}
\setlength{\parskip}{5pt}
\setlength{\emergencystretch}{3em}
\begin{document}
\section*{\texorpdfstring{Piltz Divisor Problem for k-Fold Divisor Error Terms}{Piltz Divisor Problem for k-Fold Divisor Error Terms}}\label{438}
\subsection{Definitions and mathematical statement}
For $k\ge3$, let $d_k(n)$ count ordered $k$-tuples of positive integers with product $n$. Define the error in the summatory function by
\[
 \Delta_k(x)=\sum_{n\le x}d_k(n)
 -\operatorname*{Res}_{s=1}\frac{\zeta(s)^k x^s}{s},\qquad x\ge1.
\]
The residue is the complete polynomial-in-$\log x$ main term multiplied by $x$, not just its highest-degree term. Define
\[
 \alpha_k=\inf\{a:\ \text{for every }{\varepsilon}>0,
 \ \Delta_k(x)=O_{k,{\varepsilon}}(x^{a+{\varepsilon}})\}.
\]
Determine $\alpha_k$ for every $k\ge3$. Constants may depend on $k$ and ${\varepsilon}$, but not on $x$. This is a pointwise error-exponent problem rather than a mean-square or short-interval average problem.
\par\smallskip\textit{Formulation and concept references:} \href{https://arxiv.org/html/2303.05028}{[S1]}.

\subsection{Short English statement}
What is the best possible pointwise error exponent when counting ordered factorizations into a fixed number of positive factors?
\subsection{Research attempt}
\paragraph{Process and precise target.}
This attempt was generated by GPT-6 Astra Ultra on 27 September 2026. We derive an exact multidimensional hyperbola identity, carry it through completely for $k=3$, and examine the cancellation lost in this elementary approach. The resulting exponent $2/3$ is only a baseline upper bound for $\alpha_3$, far from a determination of the exponent. No claim of a new record is made.

The pointwise exponent in the question must be distinguished from exponents defined by mean-square integrals. The introduction of [S1] records the conjectural pointwise value $(k-1)/(2k)$ and its known omega obstruction, as well as the separate mean-square problem. The sign-change results in [S2] and number-field resonance results in [S3] have their own conclusions; they do not supply the missing pointwise upper estimate here. We retain the complete residue polynomial throughout, since subtracting only its highest-degree term changes the error by an amount of order $x$ times a logarithmic power.

\paragraph{An exact identity before any estimation.}
Let $D_r(x)=\sum_{n\le x}d_r(n)$ for $r\ge1$, with $D_0(t)=1$ for $t\ge1$. Fix $k\ge2$ and $y=\lfloor x^{1/k}\rfloor$. Every ordered $k$-tuple of positive integers with product at most $x$ has at least one coordinate at most $y$: otherwise its product is at least $(y+1)^k>x$. Inclusion-exclusion on these $k$ coordinate events gives
\[
 D_k(x)=\sum_{j=1}^k(-1)^{j-1}\binom kj
 \sum_{1\le a_1,\ldots,a_j\le y}
 D_{k-j}\!\left(\frac{x}{a_1\cdots a_j}\right).
 \tag{438.1}
\]
When $j=k$, every selected tuple is admissible because its product is at most $y^k\le x$, and that term equals $(-1)^{k-1}y^k$. The identity counts ordered tuples, as required by $d_k$; it would not be valid with unordered factorizations and the same binomial coefficients.

For $k=3$ this becomes
\[
 D_3(x)=3\sum_{a\le y}D_2(x/a)
       -3\sum_{a,b\le y}\left\lfloor\frac{x}{ab}\right\rfloor+y^3.
 \tag{438.2}
\]
This exact relation is also useful for finite verification because its two sides can be computed by different counting procedures.

\paragraph{The two-dimensional input with its error accounted for.}
For $t\ge1$, set $z=\lfloor\sqrt t\rfloor$. Splitting lattice points below a hyperbola yields
\[
 D_2(t)=2\sum_{a\le z}\lfloor t/a\rfloor-z^2
       =2tH_z-z^2+O(z),
\]
where $H_z=\sum_{a\le z}1/a$. Using $H_z=\log z+\gamma+O(1/z)$, $z^2=t+O(\sqrt t)$, and $\log z=\tfrac12\log t+O(t^{-1/2})$ gives
\[
 D_2(t)=t\log t+(2\gamma-1)t+O(\sqrt t).
 \tag{438.3}
\]
The estimate is uniform for $t\ge1$ after enlarging the absolute constant over a bounded initial interval. Inserting (438.3) in the first sum of (438.2) costs
\[
 O\!\left(\sqrt x\sum_{a\le y}a^{-1/2}\right)
 =O(\sqrt{xy})=O(x^{2/3}).
\]
Replacing each floor in its second sum by its argument costs $O(y^2)=O(x^{2/3})$. Thus, with $L=\log x$ and $J_y=\sum_{a\le y}(\log a)/a$,
\[
 D_3(x)=3x\{(L+2\gamma-1)H_y-J_y-H_y^2\}+y^3+O(x^{2/3}).
 \tag{438.4}
\]
No cancellation between the floor errors has been used at this point.

\paragraph{Recovering the complete main polynomial.}
Fix the convention
\[
 \zeta(1+u)=u^{-1}+\gamma-\gamma_1u+O(u^2).
\]
In this convention $\gamma_1=\lim_{y\to\infty}(J_y-\tfrac12\log^2y)$. The first Euler summation corrections give, for integral $y\to\infty$ and $\ell=\log y$,
\[
 H_y=\ell+\gamma+\frac1{2y}+O(y^{-2}),\qquad
 J_y=\frac{\ell^2}{2}+\gamma_1+\frac{\ell}{2y}
                   +O((1+\ell)y^{-2}).
 \tag{438.5}
\]
These estimates follow by applying the trapezoidal summation formula to $1/t$ and $(\log t)/t$ on the tail; the derivatives have integrable variation of the asserted orders. The limits specify the constants and prevent a sign change in $\gamma_1$.

Here $\ell=L/3+O(1/y)$. The expression in braces in (438.4), omitting the $1/y$ corrections, is
\[
 (L+2\gamma-1)(\ell+\gamma)-\tfrac12\ell^2-\gamma_1-(\ell+\gamma)^2.
\]
Its derivative in $\ell$ is $L-3\ell-1$, bounded when $\ell=L/3+O(1/y)$. Replacing $\ell$ by $L/3$ therefore changes it by $O(1/y)$, without an extra factor $L$. Likewise, the combined coefficient of $1/(2y)$ in (438.5) is $L-3\ell-1$, again bounded. The remaining error contributes $O(x(1+L)/y^2)$, which is $O(x^{2/3})$. Finally $y^3=x+O(x^{2/3})$. Simplification yields
\[
 \begin{split}
 D_3(x)&=xP_2(\log x)+O(x^{2/3}),\\
 P_2(L)&=\tfrac12L^2+(3\gamma-1)L
              +1-3\gamma+3\gamma^2-3\gamma_1.
 \end{split} \tag{438.6}
\]
The derivation was for sufficiently large $x$; a change of constant covers $x\ge1$.

To verify that this is the main term specified in the problem, expand
\[
 \zeta(1+u)^3=u^{-3}+3\gamma u^{-2}
                       +(3\gamma^2-3\gamma_1)u^{-1}+O(1),
\]
and $x^{1+u}/(1+u)=x[1+(L-1)u+(L^2/2-L+1)u^2+O(u^3)]$. The coefficient of $u^{-1}$ in their product is exactly $xP_2(L)$. Consequently (438.6) proves $\alpha_3\le2/3$ with the catalogue's normalization, not merely a bound after an approximate subtraction.

\paragraph{Where the attempted improvement fails.}
One might try replacing the absolute bounds on the two error sums by square-root cancellation. The second is a sum of fractional parts evaluated at $x/(ab)$ over a highly dependent multiplication table. The values are neither independent random variables nor uniformly distributed by any argument above. The first contains errors $\Delta_2(x/a)$ at correlated arguments. An upper bound for each term gives the exponent already obtained; asserting random signs would be an additional unproved hypothesis.

The target suggested by the literature for $k=3$ is $1/3$, so saving a small power from a single error estimate would still not establish the required endpoint exponent. Moreover $\alpha_3$ is an infimum with an $\varepsilon$ allowance. Determining it does not require a bound with exactly $x^{1/3}$ and no logarithms, but it does require estimates for every positive $\varepsilon$, with constants independent of $x$.

Nor can a mean-square estimate simply be read as a pointwise one. If $\int_X^{2X}|E(t)|^2dt\le X^{1+2b}$, Chebyshev's inequality controls the measure of the set where $|E(t)|>X^{b+\eta}$ by $X^{1-2\eta}$. It does not prove that this exceptional set is empty. Narrow spikes are compatible with such an integral bound. The actual divisor error has arithmetic structure, but a separate argument exploiting that structure is necessary to eliminate every large exceptional value.

For larger $k$, (438.1) remains available. Its usefulness is organizational: it records intersections exactly and makes the family of correlated lower-dimensional errors explicit. It does not automatically determine their joint cancellation. Treating those errors as independent would import the conclusion rather than derive it.

\paragraph{Verification and outcome.}
The finite checks compare (438.1) with divisor counts computed by convolution for several $k$ and all integer $x$ in a fixed initial range. Symbolic coefficient arithmetic separately checks the residue polynomial and the cancellations of the apparent $x\log x/y$ terms. These calculations test identities, not the asymptotic optimal exponent.

\textbf{UNRESOLVED.} The exact hyperbola decomposition and the full $k=3$ calculation give a checked elementary upper bound. No optimal pointwise exponent is established for any of the general $k\ge3$ targets, and the missing cancellation estimate is explicitly unproved.

\subsection{Sources}
\begin{itemize}
\item[S1]Bellotti and Yang, On the generalised Dirichlet divisor problem (2023 preprint; BLMS 2024).
\url{https://arxiv.org/html/2303.05028}.
\textit{Formulation source.}
\item[S2]Status source for Piltz Divisor Problem for k-Fold Divisor Error Terms.
\url{https://academic.oup.com/imrn/article/2024/20/13191/7755133}.
\textit{Further reference; not a proof endorsement.}
\item[S3]Status source for Piltz Divisor Problem for k-Fold Divisor Error Terms.
\url{https://arxiv.org/html/2506.22587v1}.
\textit{Further reference; not a proof endorsement.}
\end{itemize}
\end{document}
